\section{Instantiation}\label{sec:txid-instantiation}

We describe a concrete instantiation of the black-box construction in pairing
based groups using algebraic PRFs and tailored sigma protocols, yielding a
transparent setup. Three of the four general relations involve pairing equations
that would result in computationally expensive proofs. We instead rewrite the
equations to require only proofs of statements involving discrete logarithms, so
that they can be handled by efficient sigma protocols. In each case the verifier
still evaluates one or two pairings itself, on public data.  What the rewriting
removes is the need to prove anything about a pairing.  $\rel_6$ and $\rel_7$
are rewritten as $\rel_6'$ and $\rel_7'$ below; the $\group{2}$ clause of
$\rel_5$ is rearranged in place and then batched.

The setup is transparent in the strong sense that no party, including the
registration authority, samples the group generators. All of
$\gen_0,\ldots,\gen_4 \in \group{1}$ and $\ggen,\ggen_D \in \group{2}$ are
derived by hash-to-curve from a public seed fixed in $\crs$, so no discrete
logarithm relation among them is known to anyone. This matters because the same
generators serve as the registration authority's BBS+ public key and as the
Pedersen commitment key, and Pedersen binding fails outright against a party
that chose the generators. The registration authority samples only its signing
scalar $z$.

\subsection{Instantiating $\FE$}

The tag PRF used below takes two inputs of different kinds: a $\group{1}$
element, which will be the epoch key, and a $\Zp$ scalar, which will be the key
it inverts against. The tracing key is accordingly a pair of field elements,
$\utk = (\utk', \utk'') \in \Zp^2$. Here $\utk'$ is the exponent the epoch PRF
$\FE$ raises to, producing the $\group{1}$ element, and $\utk''$ is the scalar
the tag PRF $\FT$ inverts against. Both components are drawn once at
registration and certified together in the same credential, and $\utk''$ is
unchanged across epochs.

\begin{assumptionbox}{Decisional Diffie-Hellman (DDH)}{ddh-txid}

  The DDH assumption holds in $\group{1}$ if no PPT adversary can distinguish
$(\gen^x, \gen^y, \gen^{xy})$ from $(\gen^x, \gen^y, \gen^z)$ for uniformly
random $x, y, z \sample \Zp$.

\end{assumptionbox}

We instantiate $\FE$ by fixing the generator to $\HE(\epoch)$, for hash function
$\HE : \{0,1\}^* \to \group{1}$ modelled as a random oracle, and the PRF key to
$\utk'$. The epoch tracing key is then the pair given below:
%
\begin{equation*}
%
  \ek = \FE(\utk, \epoch) = \big(\ek',\ \{W_\gamma\}_{\gamma\in[0,\Max)}\big),
\qquad \ek' = \HE(\epoch)^{\utk'}, \qquad W_\gamma =
\ggen_D^{1/(\utk''+\gamma)}.
%
\end{equation*}
%
This is the split of~\Cref{eq:fe-split}, with $\ek_\epoch = \ek'$ and
$\mathsf{aux} = \{W_\gamma\}_{\gamma\in[0,\Max)}$. Only $\ek'$ varies with the
epoch. The group elements $\{W_\gamma\}$ are static across epochs, and are the
outputs of the tag PRF's inner Dodis-Yampolskiy evaluation. We return below to
why the auditor receives them rather than $\utk''$ itself.
\Cref{lem:txid-fe-prf} discharges condition (F1) and \Cref{lem:txid-ft-prf}
discharges (F2); note that $\FE$ itself is emphatically \emph{not} a
pseudorandom function, since $\mathsf{aux}$ is the same in every epoch, which is
precisely why~\Cref{sec:txid-construction} states the requirement as (F1) and
(F2) rather than as pseudorandomness of $\FE$.

\begin{lemma}\label{lem:txid-fe-prf}
  %
  Under the DDH assumption in $\group{1}$ and modelling $\HE$ as a random
oracle, the epoch-varying component $\epoch \mapsto \HE(\epoch)^{\utk'}$ is a
secure pseudorandom function. This is condition (F1)
of~\Cref{sec:txid-construction}.
%
\end{lemma}

\subsection{Instantiating $\FT$}

\begin{assumptionbox}{Multi-instance Decisional Bilinear Diffie-Hellman
Inversion with published elements ($(q,n)$-mDBDHI)}{qmdbdhi-txid}

  Let $(\group{1},\group{2},\group{T},\pairing)$ be a Type-3 bilinear group of
prime order $p$ with generators $\gen \in \group{1}$ and $\ggen_D \in
\group{2}$. The $(q,n)$-mDBDHI assumption holds if, given $(\gen, \ggen_D,
\{\ggen_D^{1/(x+\gamma)}\}_{\gamma\in[0,q)})$ for uniformly random $x \sample
\Zp$ and $A_1,\ldots,A_n \sample \group{1}$, no PPT adversary can distinguish
$\{\pairing(A_i, \ggen_D^{1/(x+\gamma)})\}_{\gamma\in[0,q),\, i\in[1,n]}$ from
$\{\Delta_{i,\gamma}\}_{\gamma \in [0,q),\, i\in[1,n]}$ for uniformly random
$\Delta_{i,\gamma} \sample \group{T}$ with non-negligible advantage.  Neither
any $A_i$ nor $x$ is given to the adversary. We use it with $q = \Max$ and $n$
the number of epochs.

\end{assumptionbox}

The $n$ instances share the base elements $\{\ggen_D^{1/(x+\gamma)}\}$, which is
what the protocol demands: $\{W_\gamma\}$ is fixed at registration and reused in
every epoch, while the epoch key varies. The multi-instance version reduces to
the single-instance one ($n=1$) at no loss, by random self-reducibility: given a
single challenge $\{\Delta_\gamma\}$, set $\Delta_{i,\gamma} \leftarrow
\Delta_\gamma^{a_i}$ for fresh $a_i \sample \Zp^*$. In the real case this is a
valid $n$-instance challenge with $A_i = A^{a_i}$, and in the random case each
row is uniform in $\group{T}$, so the two distributions map to the two
distributions. We therefore state and use the multi-instance form directly and
do not carry a hybrid factor.

This is the decisional variant of the $q$-BDHI assumption~\cite{EC:BonBoy04},
with the base elements shared and the $\group{1}$ element withheld. Similar
$q$-type assumptions are the norm for constructions of this shape: a comparable
partially oblivious PRF built on the Dodis-Yampolskiy PRF with a secret group
element is analysed under a strong Diffie-Hellman inversion assumption of its
own~\cite{EC:TCRSTW22}, and it is known that pairing-based constructions of this
type cannot have both short proofs and a reduction to a simple
assumption~\cite{TCC:BHKU22}.

\begin{primitivebox}{The Dodis-Yampolskiy pseudorandom function}{fig:dy-prf}

  \begin{description}

    \item[$\ggen_D \sample \PRFgen(\secparam)$ :] On input the security
parameter, sample and return a generator $\ggen_D \sample \group{2}$.

    \item[$y \leftarrow \PRFeval(\ggen_D, X, k, x)$ :] On input generator
$\ggen_D$, group element $X \in \group{1}$, key $k \in \Zp$, and input $x \in
\Zp \setminus \{-k\}$, return the following:
      %
      \begin{equation*}
      %
	y = \pairing\big(X, \ggen_D^{1/(k+x)}\big).
      %
      \end{equation*}

  \end{description}

\end{primitivebox}

The Dodis-Yampolskiy PRF is pseudorandom under the $q$-BDHI
assumption~\cite{PKC:DodYam05}.

We instantiate $\FT$ by fixing $X$ to the epoch tracing key $\ek'$ rather than a
fixed public generator, and the PRF key $k$ to $\utk''$, giving $\FT(\ek,\gamma)
= \PRFeval(\ggen_D, \ek', \utk'', \gamma) = \pairing(\ek', W_\gamma)$ for the
$W_\gamma$ defined above.

We use the Dodis-Yampolskiy PRF as a black box, so its key $\utk''$ must stay
secret for pseudorandomness to hold. We therefore give the auditor its outputs
rather than its key. The auditor receives $\{W_\gamma\}$ as part of $\ek$, and
tracing is unaffected, since an auditor holding $\ek$ still computes every tag
of the epoch directly as $\pairing(\ek', W_\gamma)$.

Releasing $\utk''$ instead would break epoch binding outright. Writing $Y =
\pairing(\ek',\ggen_D)$, the tags of an epoch are $\Tag_\gamma =
Y^{1/(\utk''+\gamma)}$. An auditor holding $\utk''$ recovers $Y =
\Tag_\gamma^{\utk''+\gamma}$ from any single tag on the ledger, and from $Y$
computes every other tag of that user in that epoch. Since $\utk''$ is the same
in every epoch, one grant would let that auditor trace the user in all later
epochs, which is exactly the property the epoch key is meant to deny them.
Publishing $\{W_\gamma\}$ removes the auditor's need for the scalar, and
$\utk''$ never leaves the user. The epoch key $\ek'$ blinds the released
outputs, and the tags rest on that blinding.

Because $\{W_\gamma\}$ is fixed at registration, every auditor a user has ever
granted holds the same copy of it, and it is therefore a persistent pseudonymous
identifier for that user. This leaks nothing in our model: permissions are
indexed by $(\auditor,\user,\epoch)$, so an auditor is told which user it is
being authorised to trace, and two colluding auditors who compare their copies
of $\{W_\gamma\}$ learn only that they were granted by the same user, which each
of them already knew. What they cannot do is combine their copies with a tag
from an epoch for which neither holds $\ek'$, and that is exactly the statement
of~\Cref{lem:txid-ft-prf}.

\begin{lemma}\label{lem:txid-ft-prf} Let $\ek'_1,\ldots,\ek'_n$ be uniform and
independent in $\group{1}$ and unknown to the adversary, and let
$\{W_\gamma\}_{\gamma\in[0,\Max)}$ be published. Under the $(\Max,n)$-mDBDHI
assumption, the tag vectors $\FT(\ek_i,\gamma) = \pairing(\ek'_i, W_\gamma)$ are
jointly indistinguishable from independent uniform elements of $\group{T}$; in
particular $\FT$ is a secure pseudorandom function on domain $[0,\Max)$ in each
epoch. This is condition (F2) of~\Cref{sec:txid-construction}.  \end{lemma}

\subsection{Instantiating the tracing scheme}

The BBS+ signature scheme enables signers to sign multiple messages at once,
using the following
algorithms~\cite{C:BonBoySha04,SCN:AuSusMu06,EPRINT:CamDriLeh16}.

\begin{primitivebox}{The BBS+ signature scheme}{fig:bbs}

  \begin{description}

    \item[$(\sk,\pk) \sample \BBSgen(\secparam, n, (\gen_0,\ldots,\gen_{n+1},
\ggen))$ :] On input the security parameter, vector length $n$, and generators
$\gen_0,\gen_1, \ldots, \gen_{n+1} \in \group{1}$, $\ggen \in \group{2}$, do the
following:

      \begin{enumerate}

	\item Sample $z \sample \Zp$.

	\item Return $\sk = z$, $\pk = (\gen_0, \ldots, \gen_{n+1}, \ggen,
\ggen_Z = \ggen^z)$.

      \end{enumerate}

      The generators are an input rather than an output. Standard presentations
have the signer sample them, which is sound when the signer is the only party
relying on them, but here the same generators also key a Pedersen commitment
whose binding must hold \emph{against} the signer. They are therefore fixed once
in $\crs$ by hash-to-curve from a public seed, and no party knows a discrete
logarithm relation among them.

    \item[$\signature \sample \BBSsign(\sk,(\msg_1, \ldots, \msg_n))$ :] On
input secret key $\sk = z$ and message $(\msg_1, \ldots, \msg_n) \in
\group{1}^n$, do:

      \begin{enumerate}

	\item Sample $e,r \sample \Zp$.

	\item Return $\signature = (e, r, E)$ with $E = (\gen_0 \cdot
\prod_{i=1}^{n} \msg_i \cdot \gen_{n+1}^{r})^{1/(z+e)}$.

      \end{enumerate}

    \item[$\{0,1\} \leftarrow \BBSverify(\pk,(\msg_1, \ldots, \msg_n),
\signature)$ :] On input public key $\pk$, message $(\msg_1, \ldots, \msg_n)$,
and signature $\signature = (e,r,E)$, do:

      \begin{enumerate}

	\item Return $\pairing(E,\ggen_Z\cdot\ggen^{e}) \checkcheck
\pairing(\gen_0 \cdot \prod_{i=1}^{n} \msg_i \cdot \gen_{n+1}^{r}, \ggen)$.

      \end{enumerate}

  \end{description}

\end{primitivebox}

BBS+ signatures are existentially unforgeable under the $q$-strong
Diffie-Hellman ($q$-SDH) assumption.

\begin{assumptionbox}{$q$-Strong Diffie-Hellman ($q$-SDH)}{qsdh-bbs}

  The $q$-SDH assumption holds in the Type-3 bilinear group
$(\group{1},\group{2},\group{T},\pairing)$ if, given $(\gen_0, \gen_0^{\theta},
\ldots, \gen_0^{\theta^{q}}, \ggen, \ggen^{\theta})$ for $\gen_0 \in \group{1}$,
$\ggen \in \group{2}$ and uniformly random $\theta \sample \Zp$, no PPT
adversary can output a pair $\big(\varsigma,\,
\gen_0^{1/(\theta+\varsigma)}\big)$ for any $\varsigma \in \Zp \setminus
\{-\theta\}$ with non-negligible probability.

\end{assumptionbox}

This variant of the BBS+ signature scheme signs the group elements $M_i =
\gen_i^{m_i}$ directly, rather than the scalars $m_i$ themselves. Because
computing $\prod_i M_i$ requires no knowledge of the discrete logarithms $m_i$,
a signer can produce a valid signature on committed values without ever learning
them. This is exactly what lets the registration authority certify a user's keys
below without ever seeing them~\cite{SCN:AuSusMu06,EC:TesZhu23b}.

Unforgeability of this variant carries a condition that the scalar-message
version does not. An adversary who could obtain a signature on an arbitrary
group element could request one on $M_i' = M_i \cdot \gen_i^{\delta}$ and shift
the certified message by a quantity it knows, so the reduction to $q$-SDH
requires that the signer only ever signs elements accompanied by a proof of
knowledge of their representation in the bases $\gen_i$. The registration
protocol below supplies exactly that, in the shape of $\algoproof_4$, and
$\rel_4$ is therefore not an optional hardening step but a precondition of the
signature scheme's security in this setting~\cite{SCN:AuSusMu06,EC:TesZhu23b}.

\begin{primitivebox}{The Pedersen commitment scheme}{fig:pedersen}

  \begin{description}

    \item[$\ck \sample \ComGen(\secparam, n)$ :] On input the security parameter
and message vector length $n$, sample $\gen_1,\ldots,\gen_n, \hgen \sample
\group{1}$ and return $\ck = (\gen_1,\ldots,\gen_n,\hgen)$.

    \item[$\com \leftarrow \ComCommit(\ck,(m_1,\ldots,m_n),\openinfo)$ :] On
input commitment key $\ck$, messages $(m_1,\ldots,m_n) \in \Zp^n$, and opening
$\openinfo \in \Zp$, return $\com = \prod_{i=1}^{n} \gen_i^{m_i}\cdot
\hgen^{\openinfo}$.

    \item[$\{0,1\} \leftarrow \ComOpen(\ck,\com,(m_1,\ldots,m_n),\openinfo)$ :]
On input $\ck$, commitment $\com$, messages $(m_1,\ldots,m_n)$, and opening
$\openinfo$, return $1$ if $\com = \prod_{i=1}^{n}
\gen_i^{m_i}\cdot\hgen^{\openinfo}$, and $0$ otherwise.

  \end{description}

\end{primitivebox}

\paragraph{The Pedersen commitment scheme.} This scheme is perfectly hiding and
computationally binding under the discrete log assumption. Commitments are group
elements.

\begin{protocolbox}{Tracing protocol}{prot:txid}

  \begin{description}

    \item[$(\crs,\state) \sample \IDgen(\secparam)$ :] On input the security
parameter, do the following:

      \begin{enumerate}

	\item Fix bilinear groups $(\group{1},\group{2},\group{T},p,\pairing)$
of prime order $p$ admitting an efficient pairing $\pairing: \group{1} \times
\group{2} \to \group{T}$.

	\item Fix a public seed $\mathsf{sd}$ and derive
$\gen_0,\gen_1,\gen_2,\gen_3,\gen_4 \in \group{1}$ and $\ggen,\ggen_D \in
\group{2}$ by hash-to-curve from $\mathsf{sd}$.  No party samples these, so no
discrete logarithm relation among them is known to anyone, including the
registration authority. Fix a hash function $\HE : \{0,1\}^* \to \group{1}$,
modelled as a random oracle, and a hash function $\hash : \{0,1\}^* \to \Zp$.

	\item The registration authority samples $(\sk_{RA},\pk_{RA}) \sample
\BBSgen(\secparam,3,(\gen_0,\ldots,\gen_4,\ggen))$, that is, it samples only the
scalar $z$ and publishes $\ggen_Z = \ggen^z$.

	\item Set $\ck \leftarrow (\gen_2,\gen_3,\gen_4)$, reusing the
generators of Step 2, so that $(\gen_2,\gen_3)$ commit to $(\utk',\utk'')$ and
$\gen_4$ plays the role of the Pedersen hiding generator.

	\item Set $\Max \leftarrow 2^{\ell}$, $\FE(\utk,\epoch) := (\ek',
\{W_\gamma\}_{\gamma\in[0,\Max)})$ for $\ek' = \HE(\epoch)^{\utk'}$ and
$W_\gamma = \ggen_D^{1/(\utk''+\gamma)}$, and $\FT(\ek,\gamma) := \pairing(\ek',
W_\gamma)$, following the Dodis-Yampolskiy construction~\cite{PKC:DodYam05}.

	\item Set $\crs \leftarrow (\group{1},\group{2},\group{T},\pairing,
\mathsf{sd}, \gen_0,\gen_1,\gen_2,\gen_3,\gen_4,\ggen,\ggen_D,\HE,\hash,
\pk_{RA},\ck,\Max)$ and $\state \leftarrow \varnothing$, the empty ledger,
indexed by tag.

	\item Return $(\crs,\state)$.

      \end{enumerate}

    \item[$(\sk,\regcred) \sample \IDregister(\crs,\state)$ :] Interactively
between the user and the registration authority, do:

      \begin{enumerate}

	\item The registration authority sends a fresh nonce $\rand$ to the
user.

	\item The user samples $\sk,\utk',\utk'' \sample \Zp$, sets $S
\leftarrow \gen_1^{\sk}$ and $T \leftarrow \gen_2^{\utk'}\cdot\gen_3^{\utk''}$,
computes $\algoproof_4 \sample \NIZKprove(\crs_4,\inst_4,\witness_4,\rand)$ for
$\inst_4 = (\gen_1,\gen_2,\gen_3,S,T)$ and $\witness_4 = (\sk,\utk',\utk'')$,
proving $S = \gen_1^{\sk} \wedge T = \gen_2^{\utk'}\cdot\gen_3^{\utk''}$, and
sends $(S,T,\algoproof_4)$ to the registration authority. Here $\algoproof_4$ is
a standard Schnorr proof of representation over the bases
$(\gen_1,\gen_2,\gen_3)$, made non-interactive by the Fiat-Shamir transform with
$\rand$ included in the challenge hash.

	\item If $\algoproof_4$ does not verify, or the user has registered
before, the registration authority exits.

	\item Else, the registration authority computes $\signature \sample
\BBSsign(\sk_{RA},(S,T))$, that is, $\signature = (e,r,E)$ with $E =
(\gen_0\cdot S\cdot T\cdot\gen_4^{r})^{1/(z+e)}$, and records $(S,T)$ as
registered. The three message slots of the credential are
$(\gen_1^{\sk},\gen_2^{\utk'},\gen_3^{\utk''})$, and the last two reach the
registration authority bundled as the single element $T$.

	\item Return $(\sk,\utk',\utk'')$ as the user's secret key, and
$\regcred \leftarrow (S,T,\signature)$ as their registration credential.

      \end{enumerate}

    \item[$\ek \sample \IDgrant(\crs,\sk,\utk',\utk'',\regcred,\epoch)$ :]
Interactively between the user and an auditor, on a fresh nonce $\rand$ from the
auditor, do:

      \begin{enumerate}

	\item Parse $\regcred = (S,T,\signature)$.

	\item The user computes $\ek' \leftarrow \HE(\epoch)^{\utk'}$ and
$W_\gamma \leftarrow \ggen_D^{1/(\utk''+\gamma)}$ for each $\gamma \in
[0,\Max)$, and sets $\ek \leftarrow (\ek',\{W_\gamma\})$.

	\item Both parties derive $t_\gamma = \hash(\{W_\gamma\},\gamma) \in
\Zp$ and set $\widehat{W} = \prod_{\gamma=0}^{\Max-1} W_\gamma^{t_\gamma}$ and
$\widehat{V} = \prod_{\gamma=0}^{\Max-1} (\ggen_D\cdot
W_\gamma^{-\gamma})^{t_\gamma}$.

	\item The user samples $\alpha,\beta \sample \Zp$, computes $A =
\gen_2^{\alpha}\cdot\gen_3^{\beta}$, $B = \HE(\epoch)^{\alpha}$, and $K =
\widehat{W}^{\beta}$, sets $c = \hash(\inst_5,A,B,K,\rand)$, $\eta = \alpha +
c\cdot\utk'$ and $\zeta = \beta + c\cdot\utk''$, and forms $\algoproof_5
\leftarrow (A,B,K,\eta,\zeta)$, proving $\rel_5$.

	\item The user computes $\algoproof_4 \sample
\NIZKprove(\crs_4,\inst_4,\witness_4,\algoproof_5\|\rand)$, re-authenticating
$(S,T)$ on the fresh nonce, and sends
$(S,T,\signature,\ek,\algoproof_5,\algoproof_4)$ to the auditor.

	\item The auditor accepts and outputs $\ek$ if and only if
$\algoproof_4$ verifies, $\BBSverify(\pk_{RA},(S,T),\signature) = 1$, and
$\gen_2^{\eta}\cdot\gen_3^{\zeta} \checkcheck A\cdot T^{c}$, $\HE(\epoch)^{\eta}
\checkcheck B\cdot\ek'^{c}$, and $\widehat{W}^{\zeta} \checkcheck
K\cdot\widehat{V}^{c}$.

      \end{enumerate}

    \item[$\tx \sample
\IDsubmit(\crs,\sk,\utk',\utk'',\regcred,\epoch,\payload)$ :] On input
$(\crs,\sk,\utk',\utk'',\regcred,\epoch,\payload)$ from the user, do:

      \begin{enumerate}

	\item Parse $\regcred = (S,T,\signature)$ and compute $\ek =
\FE(\utk,\epoch) = (\ek',\{W_\gamma\})$.

	\item Take a fresh, unused index $0 \leq \gamma < \Max$ from the user's
local pool for this epoch, and compute $\Tag_\gamma \leftarrow \FT(\ek,\gamma) =
\pairing(\ek',W_\gamma)$. Indices need not be used in order, so this step
requires no coordination between concurrent submissions beyond the allocation
itself.

	\item Sample $\openinfo, r_\gamma \sample \Zp$ and compute $\com \sample
\ComCommit(\ck,(\utk',\utk''),\openinfo) = \gen_2^{\utk'} \cdot \gen_3^{\utk''}
\cdot \gen_4^{\openinfo}$ together with the index commitment $\com_\gamma =
\gen_3^{\gamma}\cdot\gen_4^{r_\gamma}$.

	\item Compute $\algoproof_7' \sample
\NIZKprove(\crs_7,\inst_7',\witness_7')$, proving relation $\rel_7'$,
following~\Cref{sec:tag-computation}, together with the Bulletproofs range
argument $\algoproof_{\mathsf{rng}}$ that $\com_\gamma$ opens to a value in
$[0,\Max)$.

	\item Compute $\algoproof_6' \sample
\NIZKprove(\crs_6,\inst_6',\witness_6',\msg)$ for $\msg =
\payload\|\Tag_\gamma\|\com\|\com_\gamma\|\algoproof_7'\|
\algoproof_{\mathsf{rng}}$, proving relation $\rel_6'$,
following~\Cref{sec:tx-auth}. Everything else in the transaction is bound into
$\msg$, so no component is malleable in isolation.

	\item Set $\tx \leftarrow (\payload,\Tag_\gamma,\com,\com_\gamma,
\algoproof_7',\algoproof_{\mathsf{rng}},\algoproof_6')$ and submit $\tx$ to the
ledger.

	\item Return $\tx$.

      \end{enumerate}

    \item[$\{0,1\} \leftarrow \algo{Verify}(\crs,\state,\tx)$ :] Run by the
transaction validators. Parse $\tx$ and return $1$ if and only if all four of
the following hold:

      \begin{enumerate}

	\item $\Tag_\gamma$ does not appear in $\state$, that is, it has not
accompanied a previously accepted transaction. This is what stops a user reusing
an index within an epoch.

	\item $\algoproof_7'$ verifies for $\rel_7'$ on $\inst_7'$.

	\item $\algoproof_{\mathsf{rng}}$ verifies for $\com_\gamma$ against the
range $[0,\Max)$. This is what stops a user producing an untraceable transaction
by tagging it with an out-of-range index.

	\item $\algoproof_6'$ verifies for $\rel_6'$ on $\msg$.

      \end{enumerate}

      On accepting, the validators append $\tx$ to $\state$, recording
$\payload$ against the key $\Tag_\gamma$.

    \item[$\{\payload\} \leftarrow \IDidentify(\crs,\ek,\epoch)$ :] On input
$(\crs,\ek,\epoch)$, do:

      \begin{enumerate}

	\item Parse $\ek = (\ek',\{W_\gamma\})$. For each $0 \leq \gamma <
\Max$, compute $\Tag_\gamma \leftarrow \pairing(\ek',W_\gamma)$ and retrieve
$\state[\Tag_\gamma]$ if present.

	\item Return the set of retrieved payloads.

      \end{enumerate}

  \end{description}

\end{protocolbox}

The protocol instantiation is described in full in~\Cref{prot:txid}. The
remaining algorithms follow~\Cref{sec:txid-construction} directly. We highlight
below only the design choices specific to this instantiation.

The registration authority's key pair is for the BBS+ signature scheme.  During
registration, the user accordingly sends only the group elements $S =
\gen_1^{\sk}$ and $T = \gen_2^{\utk'}\cdot\gen_3^{\utk''}$, together with a
proof of knowledge of the underlying secrets, rather than the secrets
themselves. The user's tracing key is the pair $(\utk',\utk'')$: $\utk'$ keys
the epoch PRF $\FE$, while $\utk''$ stays fixed across all epochs and keys the
Dodis-Yampolskiy term inside $\FT$. Both components sit inside the single
element $T$, so one credential certifies them together and one opening of $T$
ties every later proof to the same pair.

The commitment key reuses the triple of generators $(\gen_2,\gen_3,\gen_4)$ that
appear in the registration authority's public key, purely for compactness.
Randomness used at signing time and at commitment time remains independently
sampled at each use. Reusing these generators is safe because both primitives
rely on the same underlying property: nobody knows the discrete logarithm of one
relative to the other.

This is why Step 2 of $\IDgen$ derives the generators by hash-to-curve from a
public seed instead of having a party sample them. Had the registration
authority generated them, as a textbook reading of $\BBSgen$ would suggest, it
could set $\gen_3 = \gen_2^{w}$ for known $w$ and thereafter open any $\com$ to
a tracing key of its choosing, since Pedersen binding holds only against a party
ignorant of that relation. Since the commitment must bind precisely against the
registration authority, the relation must be unknown to it, and hash-to-curve
from a public seed is what makes that so without any trusted setup. Given that,
a commitment under the Pedersen commitment scheme is a discrete log statement
with no pairing, while every check in the BBS+ signature scheme on these
generators goes through a pairing, so nothing a party learns from one primitive
helps break the other.

Deriving the epoch tracing key deterministically from the long-term tracing key
and the epoch identifier is what confines an auditor's tracing power to a single
epoch, since a key derived for one epoch is useless for any other. Each
transaction tag is computed by evaluating a Dodis-Yampolskiy pairing-based PRF
at the user's epoch key and a fresh index. The tag carries two separate proofs,
one attesting that it was correctly derived from a properly committed tracing
key, and one that the committed key was signed by the registration authority.

The elements $\{W_\gamma\}$ do not depend on the epoch. They are therefore sent
once, at the user's first grant, and reused for every later epoch and every
auditor. A grant after the first carries only $\ek'$ and a constant size proof:
since $\widehat{W}$ and $\widehat{V}$ depend only on $\{W_\gamma\}$, the third
verification equation of $\rel_5$ need only be checked at the first grant, and
subsequent grants drop $K$ and the vector $\{W_\gamma\}$, leaving the two
$\group{1}$ equations against the same certified $T$. The binding of $T$ ties
both proofs to one pair $(\utk',\utk'')$.

Relation $\rel_5$ is accordingly defined for $\inst_5 = (\HE, \gen_2, \gen_3,
\ggen_D, T, \epoch, \ek', \{W_\gamma\}_{\gamma\in[0,\Max)})$ and $\witness_5 =
(\utk',\utk'')$, holding exactly when the following does:
%
\begin{gather*}
%
  T = \gen_2^{\utk'}\cdot\gen_3^{\utk''} \quad \wedge \quad \ek' =
\HE(\epoch)^{\utk'} \\
%
  \wedge \quad \forall\, \gamma \in [0,\Max): \ W_\gamma^{\utk''} = \ggen_D
\cdot W_\gamma^{-\gamma}.
%
\end{gather*}
%
The last equality is $W_\gamma = \ggen_D^{1/(\utk''+\gamma)}$ rearranged, so
that it becomes an equality of discrete logarithms against the committed
$\utk''$. Note that $\utk''$ is hidden in $\rel_5$ as well: the auditor no
longer computes $T\cdot\gen_3^{-\utk''}$ itself, so we prove the opening of $T$
instead. We prove the $\Max$ equalities in batch, with the prover showing
$\widehat{W}^{\utk''} = \widehat{V}$ in place of the $\Max$ separate equalities.
The coefficients $t_\gamma$ are drawn from all of $\Zp$ via $\hash$, and are
derived from $\{W_\gamma\}$, so they are fixed only after the elements they
check are. By the small exponents argument~\cite{EC:BelGarRab98}, if any one
element is malformed then the batched equality holds with probability at most
$1/p$.

\subsubsection{Verifiable tag computation}\label{sec:tag-computation} The user
first computes the commitment $\com =
\gen_2^{\utk'}\cdot\gen_3^{\utk''}\cdot\gen_4^{\openinfo}$ and the tag
$\Tag_\gamma = \FT(\ek, \gamma) = \pairing(\ek', W_\gamma)$, where the
$W_\gamma$ are the group elements released with the epoch key. Under the
$(\Max,n)$-mDBDHI assumption (\Cref{qmdbdhi-txid}), $\FT$ is a PRF even with
$\{W_\gamma\}$ published, which is~\Cref{lem:txid-ft-prf}.  Note that $\utk''$
is a hidden witness of $\rel_7$ below, so the auditor never needs it in order to
check a transaction.

$\rel_7$ is defined for $\crs_7 \leftarrow \NIZKgen(\rel_7)$, with $\crs_7$
describing a hash function $\hash \colon \{0,1\}^* \rightarrow \Zp$, used to
derive the verifier's challenge non-interactively via the Fiat-Shamir transform,
as seen later in the sigma protocol of~\Cref{prot:txid-tag-sigma}.  The
instance, witness, and relation are given as follows:
%
\begin{equation} \begin{gathered}
%
  \inst_7 = (\HE, \ggen_D, \gen_2, \gen_3, \gen_4, \com, \Max, \epoch,
\Tag_\gamma), \qquad \witness_7 = (\utk', \utk'', \openinfo, \gamma), \\
%
  \rel_7 = \{\ \com = \gen_2^{\utk'}\cdot\gen_3^{\utk''}\cdot\gen_4^{\openinfo}
\ \wedge \ \ek' = \HE(\epoch)^{\utk'} \\ \wedge \ \Tag_{\gamma} =
\pairing\big(\ek',\, \ggen_D^{1/(\utk''+\gamma)}\big) \ \wedge \ 0 \leq \gamma <
\Max\ \}.
%
\end{gathered} \end{equation}
%

Proving $\rel_7$ directly would require an argument over a secret $\group{1}$
element ($\ek'$) paired against a term whose $\group{2}$ exponent is also
secret, which would be expensive in practice.

Instead, the user samples $s, r_\gamma \sample \Zp$ and publishes the randomised
values below, the last of which is a Pedersen commitment to the transaction
index:
%
\begin{equation*}
%
  U = \HE(\epoch)^s, \quad V = \ek'^{s}, \quad \Xi = \Tag_\gamma^{\,s}, \quad
\com_\gamma = \gen_3^{\gamma}\cdot\gen_4^{r_\gamma}.
%
\end{equation*}
%
The commitment $\com_\gamma$ is what carries $\gamma$ between the sigma protocol
and the range argument below, so that the two speak about the same value. Since
$\pairing(V,\ggen_D) = \pairing(\ek',\ggen_D)^s$ and
$\Tag_\gamma^{\utk''+\gamma} = \pairing(\ek',\ggen_D)$, we have
$\pairing(V,\ggen_D) = \Xi^{\utk''+\gamma}$. The user then proves the equivalent
relation $\rel_7'$, proved under the same $\crs_7$, with instance, witness, and
relation given as follows:
%
\begin{equation} \begin{gathered}
%
  \inst_7' = (\inst_7, U, V, \Xi, \com_\gamma), \qquad \witness_7' =
(\witness_7, s, r_\gamma), \\
%
  \rel_7' = \{\ \com = \gen_2^{\utk'}\cdot\gen_3^{\utk''}\cdot\gen_4^{\openinfo}
\ \wedge \ U = \HE(\epoch)^s \ \wedge \ V = U^{\utk'} \\ \wedge \ \Xi =
\Tag_\gamma^s \ \wedge \ \pairing(V,\ggen_D) = \Xi^{\utk''}\cdot\Xi^{\gamma} \\
\wedge \ \com_\gamma = \gen_3^{\gamma}\cdot\gen_4^{r_\gamma} \ \wedge \ 0 \leq
\gamma < \Max\ \}.
%
  \end{gathered} \end{equation}
%
  Every equality here is now an equality of discrete logarithms in $\group{1}$
or $\group{T}$ between public bases. The single pairing $\pairing(V,\ggen_D)$ is
computed by the verifier from public data.

  We prove this in two parts: a sigma protocol for the equalities of $\rel_7'$
involving $\witness_7'$, and a separate, generic range argument that $0 \leq
\gamma < \Max$.

  \begin{protocolbox}{Tag correctness sigma protocol}{prot:txid-tag-sigma}

    \begin{description}

      \item[$\algoproof_7' \sample \NIZKprove(\crs_7, \inst_7', \witness_7')$ :]
On input $\inst_7' = (\inst_7, U, V, \Xi, \com_\gamma)$ and $\witness_7' =
(\utk', \utk'', \openinfo, \gamma, s, r_\gamma)$, do:

	\begin{enumerate}

	  \item Sample $\mu_1, \mu_2, \mu_h, \mu_s, \mu_\gamma, \mu_r \sample
\Zp$, one for each witness component in the order $(\utk', \utk'', \openinfo, s,
\gamma, r_\gamma)$.

	  \item Compute $A =
\gen_2^{\mu_1}\cdot\gen_3^{\mu_2}\cdot\gen_4^{\mu_h}$, $B =
\HE(\epoch)^{\mu_s}$, $D = U^{\mu_1}$, $C = \Tag_\gamma^{\mu_s}$, $F =
\Xi^{\mu_2 + \mu_\gamma}$, and $A_\gamma =
\gen_3^{\mu_\gamma}\cdot\gen_4^{\mu_r}$.

	  \item Compute challenge $c \leftarrow \hash(\inst_7', A, B, D, C, F,
A_\gamma)$.

	  \item Compute $\theta_1 = \mu_1 + c\cdot\utk'$, $\theta_2 = \mu_2 +
c\cdot\utk''$, $\theta_h = \mu_h + c\cdot\openinfo$, $\theta_s = \mu_s + c\cdot
s$, $\theta_\gamma = \mu_\gamma + c\cdot\gamma$, and $\theta_r = \mu_r + c\cdot
r_\gamma$.

	  \item Return $\algoproof_7' = (U, V, \Xi, \com_\gamma, A, B, D, C, F,
A_\gamma, \theta_1, \theta_2, \theta_h, \theta_s, \theta_\gamma, \theta_r)$.

	\end{enumerate}

      \item[$\{0,1\} \leftarrow \NIZKverify(\crs_7, \inst_7', \algoproof_7')$ :]
On input $\inst_7'$ and proof $\algoproof_7'$, reject unless $U \neq
1_{\group{1}}$ and $V \neq 1_{\group{1}}$, compute challenge $c \leftarrow
\hash(\inst_7', A, B, D, C, F, A_\gamma)$, and return the result of the
following checks:
      %
	\begin{gather*} \gen_2^{\theta_1}\cdot\gen_3^{\theta_2}\cdot
\gen_4^{\theta_h} \checkcheck A \cdot \com^c, \quad \HE(\epoch)^{\theta_s}
\checkcheck B \cdot U^c, \quad U^{\theta_1} \checkcheck D \cdot V^c, \\
\Tag_\gamma^{\theta_s} \checkcheck C \cdot \Xi^c, \quad \Xi^{\theta_2 +
\theta_\gamma} \checkcheck F \cdot \pairing(V,\ggen_D)^c, \\
\gen_3^{\theta_\gamma}\cdot\gen_4^{\theta_r} \checkcheck A_\gamma \cdot
\com_\gamma^{c}.  \end{gather*}

\end{description}

\end{protocolbox}

Three features of this protocol are worth drawing out before the proof.

First, the fifth equation has $\Xi$ as its only base, so it binds the sum
$\utk'' + \gamma$ and nothing finer. That is all the pairing identity needs, and
the two values are separated by the equations either side of it: $\theta_2$ is
shared with the first equation, where the binding of $\com$ pins $\utk''$, and
$\theta_\gamma$ is shared with the sixth, which opens $\com_\gamma$. The
soundness argument of~\Cref{sec:rel7-equiv} makes this precise.

Second, $\theta_s$ appears in both the second and fourth equations, over bases
in $\group{1}$ and $\group{T}$ respectively. This is what proves that the same
blinding scalar $s$ was used for $U$ and for $\Xi$, and hence that $\Xi$ is a
blinding of the tag named in the instance rather than an unrelated $\group{T}$
element.

Third, the two non-degeneracy checks matter. If $s = 0$ then $U$, $V$, and $\Xi$
all collapse to identities and the pairing equality says nothing; the check $U
\neq 1_{\group{1}}$ rules this out, since $\HE(\epoch)$ generates $\group{1}$
with overwhelming probability over the random oracle. And if $\utk'' = -\gamma$
then a prover holding $\utk' = 0$ satisfies the pairing equality with any tag at
all; the check $V \neq 1_{\group{1}}$ rules this out in turn. Both are used in
the soundness argument of~\Cref{sec:txid-security}.

The range check $0 \leq \gamma < \Max$ is established using a generic
Bulletproofs range argument~\cite{SP:BBBPWM18}, run on the same commitment
$\com_\gamma$ that the sigma protocol above already opens, and carried in the
transaction as $\algoproof_{\mathsf{rng}}$. Since $\theta_\gamma$ is the
response for the value committed in $\com_\gamma$, the value shown to lie in
$[0,\Max)$ is the same $\gamma$ that the sigma protocol extracts. As noted
in~\Cref{sec:txid-construction}, this check enforces traceability rather than
being a formality: an index outside $[0,\Max)$ would yield a tag that no auditor
ever evaluates.

\subsubsection{Transaction authorisation}\label{sec:tx-auth} Each transaction
carries a proof that attests it has originated from a registered user, following
relation $\rel_6$ defined below.

Since the user is registered using a signature $\signature = (e, r, E =
({\gen_0}\cdot\gen_1^{\sk} \cdot \gen_2^{\utk'}\cdot\gen_3^{\utk''}\cdot
\gen_4^r)^{1/(z+e)})$ from the BBS+ signature scheme, $\rel_6$ is defined for
$\crs_6 \leftarrow \NIZKgen(\rel_6)$, with $\crs_6$ describing a hash function
$\hash \colon \{0,1\}^* \rightarrow \Zp$, used, as with $\rel_7$ above, to
derive the verifier's challenge non-interactively via the Fiat-Shamir transform.
The instance, witness, and relation are given as follows:
%
\begin{equation} \begin{gathered}
%
  \inst_6 = (\pk, \com) = (\gen_0, \gen_1, \gen_2, \gen_3, \gen_4, \ggen,
\ggen_Z, \com), \qquad \witness_6 = (\sk, \utk', \utk'', \rho, e, r, E), \\
%
  \rel_6 = \{\ \com = \gen_2^{\utk'}\cdot\gen_3^{\utk''}\cdot\gen_4^{\rho} \
\wedge \\ \pairing(E, \ggen^e\cdot\ggen_Z) =
\pairing({\gen_0}\cdot\gen_1^{\sk}\cdot\gen_2^{\utk'}\cdot\gen_3^{\utk''}\cdot
\gen_4^r, \ggen)\ \}.
%
\end{gathered} \end{equation}
%

The zero-knowledge proof for $\rel_6$ on message $\msg$ is obtained by rewriting
$\rel_6$ as an equivalent relation $\rel_6'$ that avoids proofs over bilinear
equations, proved under the same $\crs_6$, following the technique of Camenisch,
Drijvers, and Lehmann for proving knowledge of a BBS+
signature~\cite{EPRINT:CamDriLeh16,USENIX:HesSinSor23}, specialised to three
messages with none disclosed. Writing $b =
{\gen_0}\cdot\gen_1^{\sk}\cdot\gen_2^{\utk'}\cdot\gen_3^{\utk''}\cdot\gen_4^r$
for the base of the credential, so that $E = b^{1/(z+e)}$, the user samples
fresh $r_1 \sample \Zp^*$ and $r_2 \sample \Zp$ and sets the following:
%
\begin{equation*}
%
  E' = E^{r_1}, \qquad r_3 = 1/r_1, \qquad d = b^{r_1}\cdot\gen_4^{-r_2}, \qquad
\bar{E} = E'^{-e}\cdot d\cdot\gen_4^{r_2}, \qquad r'' = r - r_2 r_3 - \rho.
%
\end{equation*}
%
Here $E'$ is the credential randomised by $r_1$, and $\bar{E} = E'^{-e}\cdot
b^{r_1} = E'^{z}$, which is what turns the credential check into the single
public pairing equality $\pairing(E',\ggen_Z) = \pairing(\bar{E},\ggen)$ that
the verifier evaluates itself. The randomiser generator $\gen_4$ is reused
between the credential and the commitment, so $r$ and $\rho$ must still be
reconciled. The value $r''$ is what does it. The instance, witness, and relation
are given as follows:
%
\begin{equation}\label{eq:rel6prime} \begin{gathered}
%
  \inst_6' = (\pk, \com, E', \bar{E}, d), \qquad \witness_6' = (\sk, \utk',
\utk'', \rho, e, r_2, r_3, r''), \\
%
  \rel_6' = \{\ \com = \gen_2^{\utk'}\cdot\gen_3^{\utk''}\cdot\gen_4^{\rho}\
\wedge \ \pairing(E', \ggen_Z) = \pairing(\bar{E}, \ggen) \ \wedge \\ \bar{E}/d
= E'^{-e}\cdot\gen_4^{r_2} \ \wedge \ (\gen_0\cdot\com)^{-1} = \gen_1^{\sk}\cdot
d^{-r_3}\cdot\gen_4^{r''}\ \}.
%
\end{gathered} \end{equation}
%

\begin{protocolbox}{Transaction authorisation proof}{prot:txid-tx-nizk}

  \begin{description}

    \item[$\algoproof_6' \sample \NIZKprove(\crs_6, \inst_6', \witness_6',
\msg)$ :] On input $\inst_6' = (\pk, \com, E', \bar{E}, d)$, $\witness_6' =
(\sk, \utk', \utk'', \rho, e, r_2, r_3, r'')$, and message $\msg$, do:

      \begin{enumerate}

	\item Sample $\mu_1, \mu_2, \mu_\rho, \mu_e, \mu_{r_2}, \mu_\sk,
\mu_{r_3}, \mu_{r''} \sample \Zp$, one for each witness component.

	\item Compute $t_1 = \gen_2^{\mu_1}\cdot\gen_3^{\mu_2}\cdot
\gen_4^{\mu_\rho}$, $t_2 = E'^{-\mu_e}\cdot\gen_4^{\mu_{r_2}}$, and $t_3 =
\gen_1^{\mu_\sk}\cdot d^{-\mu_{r_3}}\cdot \gen_4^{\mu_{r''}}$.

	\item Compute challenge $c \leftarrow \hash(\inst_6', t_1, t_2, t_3,
\msg)$.

	\item Compute $\theta_x = \mu_x + c\cdot x$ for each witness component
$x \in (\utk', \utk'', \rho, e, r_2, \sk, r_3, r'')$, writing $\theta_1$ and
$\theta_2$ for the responses to $\utk'$ and $\utk''$ as
in~\Cref{prot:txid-tag-sigma}.

	\item Return $\algoproof_6' = (E', \bar{E}, d, t_1, t_2, t_3,
\vec{\theta})$, for $\vec{\theta} = (\theta_1, \theta_2, \theta_\rho, \theta_e,
\theta_{r_2}, \theta_\sk, \theta_{r_3}, \theta_{r''})$.

      \end{enumerate}

    \item[$\{0,1\} \leftarrow \NIZKverify(\crs_6, \inst_6', \algoproof_6',
\msg)$ :] On input instance $\inst_6'$, proof $\algoproof_6'$, and message
$\msg$, do:

      \begin{enumerate}

	\item If $E' = 1_{\group{1}}$ or $\pairing(E', \ggen_Z) \neq
\pairing(\bar{E}, \ggen)$, return $0$.

	\item Compute challenge $c \leftarrow \hash(\inst_6', t_1, t_2, t_3,
\msg)$.

	\item Return the result of the following checks:
          %
	  \begin{gather*}
\gen_2^{\theta_1}\cdot\gen_3^{\theta_2}\cdot\gen_4^{\theta_\rho} \checkcheck
t_1\cdot\com^c, \quad E'^{-\theta_e}\cdot\gen_4^{\theta_{r_2}} \checkcheck
t_2\cdot(\bar{E}/d)^c, \\ \gen_1^{\theta_\sk}\cdot d^{-\theta_{r_3}}\cdot
\gen_4^{\theta_{r''}} \checkcheck t_3\cdot\big((\gen_0\cdot\com)^{-1}\big)^{c}.
\end{gather*}

      \end{enumerate}

  \end{description}

\end{protocolbox}


\begin{theorem}\label{thm:inst} The instantiation presented
in~\Cref{sec:txid-instantiation} securely realises $\idfu{\algo{trace}}$, in the
sense of the standalone model of~\Cref{sec:txid-definitions}, in the random
oracle model for $\HE$ and $\hash$, under the Decisional Diffie-Hellman
assumption in $\group{1}$ (for condition (F1)), the $(\Max,n)$-mDBDHI assumption
for $n$ the number of epochs (for condition (F2)), the $q$-Strong Diffie-Hellman
assumption (for the BBS+ signature scheme), and the discrete logarithm
assumption in $\group{1}$ (for Pedersen binding and for soundness of the
Bulletproofs range argument).  \end{theorem}

The proof is assembled in~\Cref{sec:txid-inst-sec}, from~\Cref{thm:txid-bb}
together with the lemmas and relation equivalences of~\Cref{sec:txid-security}.
